\chapter*{Índice de notación}
\addcontentsline{toc}{chapter}{Índice de notación}
\markboth{Índice de notación}{}
Una misma letra puede designar objetos distintos en capítulos diferentes; su significado local se fija antes de usarla.
\begin{longtable}{@{}p{.23\textwidth}p{.69\textwidth}@{}}
\toprule Símbolo & Significado \\\midrule\endhead
$\Acal,\Mcal$ & Álgebras o $\sigma$-álgebras de conjuntos.\\
$\sigma(\mathcal E)$ & $\sigma$-álgebra generada por una familia $\mathcal E$.\\
$\Bcal(X)$ & $\sigma$-álgebra de Borel del espacio topológico $X$.\\
$A\triangle B$ & Diferencia simétrica $(A\setminus B)\cup(B\setminus A)$.\\
$\ind A,\chi_A$ & Función indicadora de $A$, igual a uno en $A$ y a cero fuera.\\
$\mu^*,m^*$ & Medida exterior; medida exterior de Lebesgue.\\
$m,m_d,\lambda,\lambda_d$ & Medida de Lebesgue en $\R$ o en $\R^d$.\\
$\mu\times\nu$ & Medida producto de $\mu$ y $\nu$.\\
$\Acal\otimes\Bcal$ & $\sigma$-álgebra producto generada por rectángulos medibles.\\
$E_x,E^y$ & Secciones de un conjunto $E\subset X\times Y$.\\
$f^+,f^-$ & Partes positiva y negativa de una función real.\\
$f=g\;\mu\text{-c.t.p.}$ & Igualdad fuera de un conjunto $\mu$-nulo; igualdad casi en todas partes.\\
$f_n\to f$ en medida & $\mu\{|f_n-f|>a\}\to0$ para todo $a>0$.\\
$f_n\to f$ en media & $\int|f_n-f|\,d\mu\to0$.\\
$L^p(\mu)$ & Clases de funciones medibles con $\int|f|^p<\infty$, si $1\le p<\infty$.\\
$L^\infty(\mu)$ & Clases de funciones medibles esencialmente acotadas.\\
$\|f\|_p$ & $(\int|f|^p\,d\mu)^{1/p}$ para $p<\infty$; supremo esencial para $p=\infty$.\\
$\esssup |f|$ & Menor cota superior de $|f|$ válida casi en todas partes.\\
$p,q$ & Exponentes conjugados cuando $1/p+1/q=1$.\\
$V^*$ & Espacio de funcionales lineales continuos sobre $V$.\\
$\nu^+,\nu^-$ & Partes positiva y negativa de una medida con signo.\\
$|\nu|$ & Medida de variación total de una medida con signo o compleja.\\
$\nu\ll\mu$ & Continuidad absoluta: los conjuntos $\mu$-nulos son $\nu$-nulos.\\
$\nu\perp\mu$ & Singularidad: las medidas se concentran en conjuntos medibles disjuntos.\\
$d\nu/d\mu$ & Densidad de Radon--Nikodym cuando existe.\\
$C_c(X)$ & Funciones continuas de soporte compacto.\\
$C_0(X)$ & Funciones continuas que se anulan en el infinito.\\
$\supp f$ & Clausura del conjunto donde $f$ no se anula.\\
$\ell^p(I)$ & Espacio $L^p$ para la medida de conteo sobre el conjunto de índices $I$.\\
\bottomrule
\end{longtable}
